\section{A bounded path estimate}\label{sec:patterns}

We now use the word estimate of Lemma~\ref{lem:word-decay} to rule out
short systems of paths with extensive label comparisons.  An edge of the
image graph may occur many times in these paths.  The proof keeps track of
this repetition by counting subgraphs at successive multiplicity levels.
Comparisons then recover the square of the word weight needed for
Lemma~\ref{lem:word-decay}.  We develop the multiplicity and block argument
of \cite[Section~``A bounded-pattern estimate'']{OpenAI2026ZD} for the
types constructed here.

A \emph{position} of a path means an edge traversal, so different
positions may traverse the same edge of $\Gamma$.  A comparison consists
of two disjoint intervals of positions of equal length, together with
one of the following requirements: their words agree in the given
orders, or their words are $t_1\cdots t_h$ and
$\bar t_h\cdots\bar t_1$.
The comparisons in a system use disjoint positions.  We allow cyclic
intervals on a closed path.  Throughout this section comparisons must
also satisfy
\begin{equation}\label{eq:distinct-comparison-edges}
 \text{paired positions traverse distinct underlying undirected
 edges of $\Gamma$.}
\end{equation}
Write $H$ for the sum of the path lengths and $b_0$ for the number of
positions not used in a comparison.

\begin{theorem}[Bounded path estimate]\label{thm:bounded-patterns}
There is a constant $\varepsilon>0$, depending only on the fixed
alphabet and turn weights, with the following property.  For every
integer $K_0\geq1$, real number $C\geq1$, and integer $I\geq0$, the
probability under $\mathbb P_L$ that $\Gamma$ contains a system satisfying
all the following conditions tends to zero as $n\longrightarrow\infty$:
\begin{enumerate}
 \item There are at most $K_0$ positive-length paths.  Each is closed
 and cyclically immersed, except possibly one immersed interval path
 whose initial vertex is $x_A$ or $x_B$.
 \item The total length satisfies $L\leq H\leq CL$.
 \item There are at most $I$ comparisons, satisfying
 \eqref{eq:distinct-comparison-edges}, and $b_0\leq\varepsilon H$.
\end{enumerate}
In particular, $\varepsilon$ is chosen before $K_0,C,I$.  The threshold
in $n$ and the rate of convergence may depend on these three bounds.
\end{theorem}

Fix $K_0,C,I$ for the proof.  All constants described below as bounded
may depend on these parameters and on the fixed alphabet.  All
$o(L)$ estimates are uniform over the systems and patterns satisfying
these bounds.  We always work under the conditional law $\mathbb P_L$.
Index each closed path by choosing a starting vertex.  If a comparison
crosses one of these indexing points, split both intervals of that
comparison at the corresponding position.  The number of resulting
comparisons is still bounded in terms of $I$ and $K_0$, and their union
of paired positions has not changed.

\subsection{The image graph and its multiplicity levels}

Let $\Delta\subseteq\Gamma$ be the image of a putative system.  It has
at most $CL$ edges and girth at least $L$.  We first show that it has
bounded topological complexity.  Here the cycle rank of a finite graph
is $|E|-|V|+$ the number of connected components.

\begin{lemma}\label{lem:bounded-image-rank}
For fixed $C$, every finite graph of girth at least $L$ with at most
$CL$ edges has bounded cycle rank, uniformly in $L$.
\end{lemma}

\begin{proof}
Delete vertices of degree at most one repeatedly; this does not change
the cycle rank.  Components of the remaining graph that are circles
number at most $C$, since each has at least $L$ edges.  Remove these
components, and suppress every degree-two vertex in what remains.
The resulting multigraph $D$, if nonempty, has minimum degree at least
three.  Give each edge of $D$ the positive integer length of the path
it represents.  The sum of these lengths is at most $CL$.

Suppose $D$ has $r>0$ edges.  Choose a directed edge uniformly from
the $2r$ choices, then take a non-backtracking walk with $k$ directed
edges, choosing each continuation uniformly.  The uniform law on
directed edges is stationary: a directed edge starting at $z$ has
$\deg(z)-1$ allowed predecessors, each of which chooses it with
probability $1/(\deg(z)-1)$.  Thus the expected sum of the inherited
lengths along the walk is at most $kCL/r$.  Markov's inequality puts
at least half of the probability on walks of length at most
$2kCL/r$.  Since every individual walk has probability at most
$(2r)^{-1}2^{-(k-1)}$, there are at least $r2^{k-1}$ such short walks.

Take $k=\lceil\log_2(8r)\rceil+1$.  Then $r2^{k-1}>(2r)^2$,
whereas there are at most $(2r)^2$ ordered pairs of endpoint vertices.
Two distinct short walks therefore have the same endpoints.  Expanding
their edges gives distinct reduced paths in the original graph.  Their
union contains a cycle: in a forest there is only one reduced path
between any two vertices.  This cycle has length at most $4kCL/r$,
so the girth bound implies
\[
 r\leq4C\bigl(\lceil\log_2(8r)\rceil+1\bigr).
\]
The right side grows logarithmically with $r$, and hence this inequality
bounds $r$ in terms of $C$.  Suppression preserves cycle rank, whose
contribution from $D$ is at most $r$; the removed circle components
contribute at most $C$.
\end{proof}

Mark every vertex of $\Delta$ whose degree differs from two, and also
every initial and terminal vertex of our indexed paths.  There are no
isolated vertices.  A leaf can occur only as an endpoint of the possible
interval path, because every other path turn is immersed.  There are
therefore at most two leaves.  To see explicitly that the remaining
marks are bounded, write $\beta$ for the cycle rank and $c$ for the
number of components.  The degree sum gives
\[
 \sum_{\deg(z)\geq3}(\deg(z)-2)
 =2\beta-2c+\#\{z:\deg(z)=1\}.
\]
Lemma~\ref{lem:bounded-image-rank} bounds both the number and the total
degree of vertices of degree at least three.  There are at most $2K_0$
additional path-endpoint marks.  Each component contains a path start.
Consequently $\Delta$ is the union of a bounded number of \emph{chains},
that is, paths between marked vertices with disjoint, unmarked
degree-two interiors.  A chain may have the same marked vertex at both
ends.  Choose an orientation on every chain.

Every path traverses whole chains.  Let $m_e$ be the total number of
traversals of the undirected edge $e$ by the entire path system.  This
multiplicity is constant on the edges of each chain.  Two successive
uses of the same oriented edge by a single indexed path have starting
positions separated by at least $L$: the intervening path from the
initial vertex of that edge back to itself is a nonempty reduced closed
walk and contains a cycle.  For a path of length $h$, each orientation
can therefore occur at most $h/L+1$ times.  Summing over paths and the
two orientations yields
\begin{equation}\label{eq:bounded-multiplicity}
 m_e\leq2(H/L+K_0)\leq2(C+K_0),\qquad
 M:=\lceil2(C+K_0)\rceil.
\end{equation}

An \emph{unlabelled pattern} records the finite marked chain graph,
its positive chain lengths, the side $A$ or $B$ of each component,
the indexed paths as sequences of oriented chain traversals, and the
endpoints and directions of the comparisons.  It records whether there
is an interval path and, if so, which root is its initial vertex.
Only patterns having the structure just proved and satisfying
\eqref{eq:distinct-comparison-edges} are retained.  The underlying graph
in a pattern is an image graph: a realization must map its vertices and
edges injectively into the appropriate sides of $\Gamma$.

There are boundedly many chain traversals by
\eqref{eq:bounded-multiplicity}.  All lengths and comparison endpoints
are integers between $0$ and $CL$.  Hence the number of unlabelled
patterns is polynomial in $L$.  This includes the finitely many choices
of chain graph, traversal sequences, sides, and interval designation.

For a pattern, let $m_*:=\max_e m_e\leq M$.  Its \emph{stage $j$ graph}
$\Delta_j$ consists of all edges with $m_e\geq j$ and their endpoints,
for $1\leq j\leq m_*$.  Write $V_j$ and $E_j$ for its numbers of vertices
and edges.  Let $\iota$ be $1$ if there is an interval path and $0$
otherwise, and put
\[
 k_1=\iota,\qquad k_j=0\quad(j>1).
\]
We will fix the prescribed root only when counting stage one.  This
loss of a factor $n$ pays precisely for the two possible unpaired path
endpoints in the following incidence estimate.

\begin{lemma}\label{lem:multiplicity-incidence}
The stage graphs satisfy
\begin{equation}\label{eq:stage-incidence}
 \sum_{j=1}^{m_*}E_j=H,\qquad
 \sum_{j=1}^{m_*}V_j\leq H+\iota,\qquad
 \sum_{j=1}^{m_*}(V_j-E_j-k_j)\leq0.
\end{equation}
\end{lemma}

\begin{proof}
The first identity counts an edge once at each of its $m_e$ stages.
At a vertex $z$, put $a(z)$ equal to the number of interval endpoints
located there, counting both if they coincide; then
$\sum_z a(z)=2\iota$.  Every other traversal incidence is paired,
by a path turn, with an incidence on a different half-edge.
For any particular half-edge at $z$, all but at most $a(z)$ of its
$m_e$ incidences must therefore be paired with the other half-edges.
For sufficiently large $n$, $L>2$, so $\Delta$ has no loops; it follows
that
\[
 2\max_{e\ni z}m_e\leq\sum_{e\ni z}m_e+a(z).
\]
A vertex occurs at exactly $\max_{e\ni z}m_e$ stages.  Summing this
inequality over vertices gives
$\sum_jV_j\leq\sum_e m_e+\iota=H+\iota$.  Since $\sum_jk_j=\iota$,
the last assertion follows.
\end{proof}

\subsection{One grid for all comparisons}

The role of \eqref{eq:stage-incidence} is visible in the embedding count.
For fixed labels, choosing $V_j$ vertices, prescribing $E_j$ distinct
edges, and fixing $k_j$ roots leave a factor $n^{V_j-E_j-k_j}$.
Apart from factors subexponential in $L$, the turn frequencies also
supply one word weight for each chain. Thus
\eqref{eq:stage-incidence} permits us to combine the stage estimates
without a positive net power of $n$. It remains to sum over labels
compatible with the comparisons; the precise embedding count appears
in \eqref{eq:fixed-stage-expectation} below.

For orientation, suppose two edge-disjoint chain segments are required
to carry the same word $W$. Their turn-weight contribution is $P(W)^2$,
although only one word is free to vary. The same holds for inverse
words because $P(W^{-1})=P(W)$. Lemma~\ref{lem:word-decay} then supplies
exponential decay. In the actual system, the comparison endpoints
need not line up, and chains can be traversed repeatedly. We now
subdivide the chains into long words whose boundaries nearly agree
across comparisons; the multiplicity stages will account for the
repeated traversals.

Fix an unlabelled pattern.  Give each oriented chain coordinates from
$0$ to its length, putting edge positions at half-integers.  Split a
comparison whenever either of its path intervals crosses a chain end.
The bounded number of chain traversals gives a bounded number of
resulting subcomparisons.  In chain coordinates each is a restriction
of
\begin{equation}\label{eq:comparison-affine-map}
 x\longmapsto \sigma x+a,\qquad \sigma\in\{1,-1\},\quad a\in\mathbb Z.
\end{equation}
If $\sigma=1$, the labels in chain orientation agree; if $\sigma=-1$,
they are inverse.  This follows by incorporating into the comparison
the orientation of each of its two chain traversals.

We need blocks whose grids agree approximately under every map in
\eqref{eq:comparison-affine-map}.  Choose a fixed integer $d_1\geq1$
bounding the number of offsets, padding their list with zeros when
necessary.  Put
\[
 N=\left\lfloor L^{1/(2(d_1+1))}\right\rfloor,
 \qquad R=\lfloor\sqrt N\rfloor.
\]
Partition $[0,1)^{d_1}$ into $N^{d_1}$ boxes of side $1/N$.  Among
the $N^{d_1}+1$ vectors with coordinates $\{\ell a/L\}$,
$0\leq\ell\leq N^{d_1}$, two lie in the same box.  Subtracting their
indices gives $1\leq u\leq N^{d_1}$ such that
\[
 \operatorname{dist}(ua/L,\mathbb Z)\leq1/N
 \quad\text{for every offset }a.
\]
Choose one such $u$ deterministically from the pattern and set
\begin{equation}\label{eq:block-scale}
 s:=\frac{L}{uR}.
\end{equation}
Uniformly over patterns,
\[
 \frac{L}{N^{d_1}\sqrt N}\leq s\leq\frac LR,
 \qquad
 \operatorname{dist}(a/s,\mathbb Z)\leq R/N.
\]
In particular, $s\to\infty$, $s=o(L)$, and all the distances in the
last display tend to zero.  The lower bound even gives
$s\geq L^{(d_1+3/2)/(2(d_1+1))}$ for large $L$, so all estimates
involving $L/s$ are uniform.

On each chain take the edge positions lying in grid intervals
$[ks,(k+1)s)$ contained in its coordinate interval.  These sets are
the \emph{full blocks}.  The final fragment has fewer than $s+1$
positions.  Index underlying full blocks by $i$, and write $h_i$ for
their lengths and $m_i$ for their chain multiplicities.  Thus
$h_i=s+O(1)$ and $1\leq m_i\leq M$.  Every traversal of a chain uses
the same subdivision; a block $i$ therefore has $m_i$ separate
occurrences in the path system.

Choose $\kappa=\kappa_L>0$ tending to zero, large enough to dominate
the grid displacement $R/N$ and the rounding error $10/s$.  For each
subcomparison, match two full-block occurrences if it compares
contiguous portions of these occurrences of length at least
$(1-\kappa)s$.  Record the entire substring constraint on that overlap.
For large $L$ this gives a partial matching of occurrences.  Indeed,
the matched portion occupies more than half of each block.  Two
different proposed matches incident to the same occurrence would
therefore overlap, contrary to the disjointness of all compared
positions.  The match is recorded only once, irrespective of which of
its two occurrences is considered first.

We next bound the full-block occurrences left unmatched.
Exclude the boundedly many block occurrences
within distance $3s+O(1)$ of a subcomparison endpoint or a chain end,
measuring distance along the corresponding traversal.  A remaining
block occurrence either has every position unpaired, or lies wholly
inside one subcomparison.  In the latter case its grid interval maps
under \eqref{eq:comparison-affine-map} to an interval displaced by
at most $(R/N)s$ from a grid interval on the partner chain.  Reflection
may exchange open and closed endpoints, changing the number of
half-integer positions by at most two.  The partner interval is full,
because the source block is also far from the corresponding
subcomparison ends.  The two blocks consequently overlap in at least
$(1-\kappa)s$ compared positions, and are matched.

Let $b_*$ be the number of unmatched full-block occurrences.
There are $O(L/s)$ occurrences altogether.  Each unmatched occurrence
outside the bounded exceptional collection contributes at least $s-1$
unpaired positions.  It follows that
\begin{equation}\label{eq:block-unmatched-budget}
 sb_*\leq b_0+O(s+L/s)=b_0+o(L).
\end{equation}
The constants here count subcomparison ends and chain traversals, not
individual edges.

Now forget the occurrence indices while retaining all matches and
their substring constraints.  This gives a multigraph on the underlying
blocks; call its edges \emph{links}.  A link whose two occurrences
belong to the same underlying block is a \emph{self-link}.  Isolated
blocks are included as components.  For a component $\mathcal U$, put
\begin{equation}\label{eq:link-component-parameters}
 S_{\mathcal U}=\sum_{i\in\mathcal U}m_i,\qquad
 z_{\mathcal U}=\max_{i\in\mathcal U}m_i,\qquad
 b_{\mathcal U}=\#\{\text{unmatched occurrences of blocks in }
 \mathcal U\}.
\end{equation}
In particular, $\sum_{\mathcal U}b_{\mathcal U}=b_*$.
Figure~\ref{fig:aligned-blocks} illustrates the nearly aligned grids
and the distinction between occurrences and underlying blocks.

\begin{figure}[tbp]
\centering
\begin{tikzpicture}[x=1cm,y=1cm,font=\small,
  line cap=round,line join=round,>=Latex]
  \node[anchor=west,font=\small\bfseries] at (0,4.05)
    {(a) Nearly aligned blocks};
  \node at (3.50,3.48) {$f(x)=\pm x+a$};
  \fill[black!18] (1.40,2.75) rectangle (5.35,3.05);
  \fill[black!18] (1.40,1.95) rectangle (5.35,2.25);
  \draw (1.40,2.75) rectangle (5.60,3.05);
  \draw (1.15,1.95) rectangle (5.35,2.25);
  \node[anchor=east] at (1.00,2.90) {$f(B_i)$};
  \node[anchor=east] at (1.00,2.10) {$B_j$};
  \draw[densely dashed,black!55] (1.40,1.70) -- (1.40,3.18);
  \draw[densely dashed,black!55] (5.35,1.70) -- (5.35,3.18);
  \draw[<->] (1.40,1.45) -- (5.35,1.45);
  \node[align=center] at (3.38,0.90)
    {$J$: compared portion\\$|J|\ge(1-\kappa)s,\quad\kappa=o(1)$};
  \node[font=\footnotesize,align=center] at (3.38,0.20)
    {Each full block has $s+O(1)$ positions.};

  \node[anchor=west,font=\small\bfseries] at (7.00,4.05)
    {(b) Occurrences and underlying blocks};
  \node[anchor=west,font=\footnotesize] at (7.05,3.52)
    {Two disjoint occurrence matches};
  \node[draw,rounded corners=1pt,minimum width=1.02cm,
    minimum height=0.48cm] (i1) at (8.15,2.93) {$i^{(1)}$};
  \node[draw,rounded corners=1pt,minimum width=1.02cm,
    minimum height=0.48cm] (j1) at (11.65,2.93) {$j^{(1)}$};
  \node[draw,rounded corners=1pt,minimum width=1.02cm,
    minimum height=0.48cm] (i2) at (8.15,2.18) {$i^{(2)}$};
  \node[draw,rounded corners=1pt,minimum width=1.02cm,
    minimum height=0.48cm] (k1) at (11.65,2.18) {$k^{(1)}$};
  \draw (i1) -- (j1);
  \draw (i2) -- (k1);
  \draw[->] (9.90,1.75) -- (9.90,1.15)
    node[midway,right,font=\footnotesize] {forget occurrence indices};
  \node[draw,circle,inner sep=2.5pt] (j) at (8.15,0.64) {$j$};
  \node[draw,circle,inner sep=2.5pt,fill=black!12] (i) at (9.90,0.64) {$i$};
  \node[draw,circle,inner sep=2.5pt] (k) at (11.65,0.64) {$k$};
  \draw (j) -- (i) -- (k);
  \node[font=\footnotesize] at (9.90,0.02)
    {One underlying block $i$, two links};
\end{tikzpicture}
\caption{Aligned blocks and the two levels of comparison.  In (a), away
from subcomparison ends, the image of a full grid interval is displaced
by $o(s)$ from a partner grid interval; the shaded portion $J$ records a
comparison of $s-o(s)$ positions.  Reflection reverses order and inverts
letters.  The intervals and their small end discrepancies are schematic.
In (b), $i^{(r)}$ denotes the $r$th occurrence of block $i$.  Occurrence
joins form a partial matching, but forgetting the occurrence indices
can give several links at one underlying block.}
\label{fig:aligned-blocks}
\end{figure}


\begin{lemma}\label{lem:link-word-counts}
There is a number $\rho_L\to0$, uniform over patterns, with the following
properties.  Every word assignment counted in this lemma is required
to be reduced on each block.
\begin{enumerate}
 \item Given the reduced word on one block of a link component
 $\mathcal U$, at most $\exp(\rho_Ls|\mathcal U|)$ assignments of
 words to all its blocks satisfy the link constraints.
 \item If $\mathcal U$ contains a self-link, the total number of such
 assignments is at most $\exp(\rho_Ls|\mathcal U|)$.
 \item If $\mathcal U$ has no self-link, then
 \begin{equation}\label{eq:component-multiplicity}
  2z_{\mathcal U}\leq S_{\mathcal U}+b_{\mathcal U}.
 \end{equation}
\end{enumerate}
\end{lemma}

\begin{proof}
Root a spanning tree of the link component at the block whose word is
given.  A link from a known block determines at least $(1-\kappa)s$
letters of the next block, whose total length is at most $s+1$.
There are at most $|T|^{\kappa s+1}$ choices for its remaining letters.
Multiplying along the tree proves the first assertion, even if we
ignore reducedness and all links outside the tree when making this
upper bound.

For the second assertion, start with a block carrying a self-link.
If its coordinate map is a translation, the constraint takes the form
$t_x=t_{x+a}$ on at least $(1-\kappa)s$ consecutive positions.  The
integer $a$ is nonzero: $a=0$ would pair occurrences of the same
underlying edge, violating \eqref{eq:distinct-comparison-edges}.
The equality graph on the $h_i$ positions is a forest.  In fact all its
edges have the same nonzero step, and within each residue class modulo
$|a|$ they connect successive positions in their natural order.
There are at least $(1-\kappa)s$ distinct equality edges, leaving at
most $h_i-(1-\kappa)s\leq\kappa s+1$ components.  Choosing one letter
per component gives at most $|T|^{\kappa s+1}$ choices.

A reflection self-link cannot occur with reduced block words.  Let
$J$ be the interval of positions on which the constraint reads
$t_x=\overline{t_{a-x}}$.  Its reflected image also lies in the same
block, and both have more than half the block's positions.
Their intersection is therefore a nonempty interval invariant under
reflection.  If it has a central fixed position, the letter there
would equal its inverse, contrary to the fixed-point-free involution
on $T$.  If it has no fixed position, its two central adjacent positions
are reflections of each other and have inverse letters, contrary to
reducedness.  After counting the words on the initial self-linked
block, propagate along a spanning tree as above.  Enlarging $\rho_L$
by a fixed factor proves the second assertion.

Finally, choose a block with $m_i=z_{\mathcal U}$.  If there is no
self-link, every matched occurrence of this block has its partner
among the $S_{\mathcal U}-z_{\mathcal U}$ occurrences of other blocks.
There are at least $z_{\mathcal U}-b_{\mathcal U}$ such matched
occurrences, and distinct occurrences have distinct partners.  Hence
$z_{\mathcal U}-b_{\mathcal U}\leq
S_{\mathcal U}-z_{\mathcal U}$, which is
\eqref{eq:component-multiplicity}.
\end{proof}

\subsection{Word bins and stage expectations}

The link constraints make the number of free words small.  To combine
this fact with the probability of embedding a labelled graph, we must
also retain each block's word weight.  For its word $W_i$, in the chosen
chain orientation, record the integer
\[
 d_i:=\lfloor-\log P(W_i)\rfloor.
\]
All these words are reduced.  Since the fixed collection of positive
turn weights has a positive minimum, there is a constant $a_0>0$
depending only on those weights such that, for large $L$,
\begin{equation}\label{eq:word-bin-bounds}
 \frac{\delta s}{2}\leq d_i\leq a_0s,
 \qquad
 \#\{W\in T^{h_i}:\lfloor-\log P(W)\rfloor=d_i\}
 \leq\exp(2d_i-\delta s).
\end{equation}
Here and below only positive-weight words are counted.  For completeness,
Lemma~\ref{lem:word-decay} implies
$P(W)^2\leq e^{-2\delta h_i}$ for each word, so
$d_i\geq\delta h_i-1\geq\delta s/2$ for large $s$.
The minimum turn weight gives the upper bound.  Every word in the bin
has $P(W)>e^{-d_i-1}$, and therefore the number in the bin is at most
\[
 e^{2d_i+2}\sum_{W\in T^{h_i}}P(W)^2
 \leq e^{2d_i+2-2\delta h_i}
 \leq e^{2d_i-\delta s}.
\]
These estimates are uniform because $h_i=s+O(1)$ and $s\to\infty$.

Enlarge the pattern by recording all $d_i$.  There are $O(L/s)$ full
blocks, each with $O(s)$ possible bins.  Thus the logarithm of the
number of bin choices is
$O((L/s)\log(s+2))=o(L)$.  Together with the earlier polynomial count,
the total number of enlarged patterns is
\begin{equation}\label{eq:pattern-entropy}
 \exp(o(L)).
\end{equation}
The grid, occurrence matching, and link multigraph are determined
by the unlabelled pattern, and introduce no further choices.

Fix one enlarged pattern.  A full label assignment means labels on
all edges of its abstract image graph, read inversely on reverse
edges, giving an immersion into $F$ and satisfying the comparisons
and bins.  This is a condition on strings on the abstract graph; it
does not require an embedding in the random graph.  At stage $j$,
consider the restrictions of all such full label assignments to
$\Delta_j$.  Let $Z_j$ count their injective realizations in $\Gamma$,
with the component sides prescribed by the pattern, and with the
root fixed if $k_j=1$.  Every realization of the original enlarged
pattern gives $Z_j\geq1$ at every stage.

For a fixed one of these stage label assignments, the expected number
of its realizations is at most
\begin{equation}\label{eq:fixed-stage-expectation}
 n^{V_j-E_j-k_j}
 \exp\left(-\sum_{i:m_i\geq j}d_i+o(L)\right).
\end{equation}
To prove this, let $Q_j$ be the number of active chains, that is,
chains whose edges belong to $\Delta_j$.  They have exactly
$E_j-Q_j$ unmarked internal vertices.  At such a vertex the incoming
and outgoing chain labels form a reduced turn $t,u$, and
\eqref{eq:turn-count} bounds its possible images on its prescribed side
by
\[
 np\,w(t,u)+O(1)=np\,w(t,u)(1+O(1/n)).
\]
At each remaining unfixed vertex use the upper bound $n$; a fixed root
is marked and has only one choice.  This bounds the number of possible
vertex assignments by
\[
 n^{V_j-k_j}p^{E_j-Q_j}
 \prod_{\text{internal turns of active chains}}w(t,u)
 \exp(O(L/n)).
\]
Restrict to the assignments that are injective and compatible with all
vertex types.  Each requires $E_j$ distinct edge prescriptions.
Lemma~\ref{lem:prescriptions} bounds their joint probability by
$(np)^{-E_j}\exp(o(L))$.  Multiplication leaves
$n^{V_j-E_j-k_j}p^{-Q_j}$ times the displayed turn product.
The bounded factor $p^{-Q_j}$ is $\exp(o(L))$.  Every turn weight is
at most one, so discarding all factors except internal turns of full
blocks increases the product.  The remaining product is
\[
 \prod_{i:m_i\geq j}P(W_i)
 \leq\exp\left(-\sum_{i:m_i\geq j}d_i\right),
\]
which proves \eqref{eq:fixed-stage-expectation}.  This count uses
injective realizations of the image graph; it does not replace distinct
edge prescriptions by repeated traversals.

The number of stage label assignments has the bound
\begin{equation}\label{eq:stage-label-count}
 \exp\left(o(L)+
 \sum_{\substack{\mathcal U\text{ without self-link}\\
                  z_{\mathcal U}\geq j}}
 \left(2\min_{i\in\mathcal U}d_i-\delta s\right)\right).
\end{equation}
Indeed, a link component is represented at stage $j$ exactly when
$z_{\mathcal U}\geq j$.  The possible words on its represented blocks
are projections of assignments on the entire link component.  A
projection has cardinality no larger than the set being projected.
If the component has no self-link, choose a block attaining the
minimum bin, even when that block is absent from stage $j$.
Equation~\eqref{eq:word-bin-bounds} gives at most
$\exp(2\min_i d_i-\delta s)$ choices on this block, and
Lemma~\ref{lem:link-word-counts} bounds the extensions by
$\exp(\rho_Ls|\mathcal U|)$.  If the component has a self-link,
the second assertion of that lemma gives this last factor without
any initial word choice.  Constraints between different words, or at
chain junctions, can only reduce these counts.

Summed over components, the propagation error is
$\rho_Ls\,O(L/s)=o(L)$.  There are only boundedly many chain end
fragments, containing $O(s)$ positions in total; assigning their
letters freely costs at most $|T|^{O(s)}=\exp(o(L))$.
These observations prove \eqref{eq:stage-label-count}.  Requiring
stage labels to extend as strings was essential: it permits the
minimum-bin block to be used even at stages where it is absent.

Let $B_j$ be the product of the right sides of
\eqref{eq:fixed-stage-expectation} and
\eqref{eq:stage-label-count}, choosing their uniform error bounds.
Then
\begin{equation}\label{eq:stage-markov}
 \mathbb E_L Z_j\leq B_j,
 \qquad
 \mathbb P_L(\text{the enlarged pattern is realized})
 \leq\min_{1\leq j\leq m_*}B_j.
\end{equation}
The second inequality is Markov's inequality at each individual stage.
We will now bound the minimum by multiplying the deterministic numbers
$B_j$.

\subsection{Combining the stages}

\begin{proof}[Proof of Theorem~\ref{thm:bounded-patterns}]
Lemma~\ref{lem:multiplicity-incidence} shows that the total exponent of
$n$ in $\prod_{j=1}^{m_*}B_j$ is nonpositive. For a component
$\mathcal U$ without a self-link, write
$d_{\min}=\min_{i\in\mathcal U}d_i$.  Its positive label-count term
occurs at $z_{\mathcal U}$ stages; the negative term $-d_i$ occurs at
$m_i$ stages.  Its total contribution to the logarithm, apart from
the uniform errors, is at most
\begin{align*}
 z_{\mathcal U}(2d_{\min}-\delta s)
       -\sum_{i\in\mathcal U}m_i d_i
 &\leq (2z_{\mathcal U}-S_{\mathcal U})d_{\min}
                -\delta s z_{\mathcal U}\\
 &\leq a_0s b_{\mathcal U}
                -\frac{\delta s}{2}S_{\mathcal U}.
\end{align*}
For the last inequality, if $2z_{\mathcal U}-S_{\mathcal U}\leq0$,
use $d_{\min}\geq\delta s/2$ to bound the preceding line by
$-\delta s S_{\mathcal U}/2$.  If
$2z_{\mathcal U}-S_{\mathcal U}>0$, use
$d_{\min}\leq a_0s$ and
$2z_{\mathcal U}-S_{\mathcal U}\leq b_{\mathcal U}$ from
\eqref{eq:component-multiplicity}; the remaining term satisfies
$-\delta s z_{\mathcal U}\leq-\delta s S_{\mathcal U}/2$.

For a component with a self-link there is no positive label-count
term, and \eqref{eq:word-bin-bounds} gives
\[
 -\sum_{i\in\mathcal U}m_i d_i
 \leq-\frac{\delta s}{2}S_{\mathcal U}.
\]
The full-block occurrences account for all positions except $O(s)$
positions in end fragments.  Replacing their actual lengths
$h_i=s+O(1)$ by $s$ changes the count by $O(L/s)$, because the
multiplicities are bounded.  Thus
\[
 s\sum_{\mathcal U}S_{\mathcal U}=H+O(s+L/s)=H+o(L).
\]
Equation~\eqref{eq:block-unmatched-budget} also gives
$s\sum_{\mathcal U}b_{\mathcal U}\leq b_0+o(L)$.
There are at most $M$ stages, so summing their uniform errors still
gives $o(L)$.  We conclude that
\begin{equation}\label{eq:stage-product}
 \log\prod_{j=1}^{m_*}B_j
 \leq-\frac\delta2 H+a_0b_0+o(L).
\end{equation}

Choose, once and for all,
\begin{equation}\label{eq:epsilon-choice}
 \varepsilon:=\min\left\{\frac14,
                  \frac{\delta}{16(1+a_0)}\right\}.
\end{equation}
Both $\delta$ and $a_0$ depend only on the fixed alphabet and weights,
so this choice is independent of $K_0,C,I$.  If $b_0\leq\varepsilon H$
and $H\geq L$, the uniform error in \eqref{eq:stage-product} is small
enough, for all sufficiently large $n$, that
\[
 \log\prod_{j=1}^{m_*}B_j\leq-\frac\delta4H.
\]
Some stage therefore has
\[
 B_j\leq\exp\left(-\frac{\delta H}{4m_*}\right)
       \leq\exp\left(-\frac{\delta H}{4M}\right).
\]
Equation~\eqref{eq:stage-markov} bounds the realization probability
of this enlarged pattern by the last expression.  No independence of
the stages has been used; a full realization must survive every one
of their separate first-moment bounds.  Finally, sum over the
$\exp(o(L))$ enlarged patterns from \eqref{eq:pattern-entropy} and
use $H\geq L$.  The resulting bound
\[
 \exp\left(-\frac{\delta L}{4M}+o(L)\right)
 \longrightarrow0
\]
proves the theorem, with the stated order of quantifiers.
\end{proof}
